% ------------------------------------------------------------------------
% ------------------------------------------------------------------------
% abnTeX2: Modelo de Artigo Acadêmico em conformidade com
% ABNT NBR 6022:2003: Informação e documentação - Artigo em publicação 
% periódica científica impressa - Apresentação
% ------------------------------------------------------------------------
% ------------------------------------------------------------------------

\documentclass[
	% -- opções da classe memoir --
	article,			% indica que é um artigo acadêmico
	11pt,				% tamanho da fonte
	oneside,			% para impressão apenas no verso. Oposto a twoside
	a4paper,			% tamanho do papel. 
	% -- opções da classe abntex2 --
	%chapter=TITLE,		% títulos de capítulos convertidos em letras maiúsculas
	%section=TITLE,		% títulos de seções convertidos em letras maiúsculas
	%subsection=TITLE,	% títulos de subseções convertidos em letras maiúsculas
	%subsubsection=TITLE % títulos de subsubseções convertidos em letras maiúsculas
	% -- opções do pacote babel --
	english,			% idioma adicional para hifenização
	brazil,				% o último idioma é o principal do documento
	sumario=tradicional
	]{abntex2}


% ---
% PACOTES
% ---

% ---
% Pacotes fundamentais 
% ---
\usepackage{lmodern}			% Usa a fonte Latin Modern
\usepackage[T1]{fontenc}		% Selecao de codigos de fonte.
\usepackage[utf8]{inputenc}		% Codificacao do documento (conversão automática dos acentos)
\usepackage{indentfirst}		% Indenta o primeiro parágrafo de cada seção.
\usepackage{nomencl} 			% Lista de simbolos
\usepackage{color}				% Controle das cores
\usepackage{graphicx}			% Inclusão de gráficos
\usepackage{microtype} 			% para melhorias de justificação
\usepackage{listings}
\usepackage{xcolor}
\usepackage{amsmath}
\usepackage{float}
\lstset{
  language=C++,
  basicstyle=\ttfamily\small,
  keywordstyle=\color{blue},
  stringstyle=\color{red},
  commentstyle=\color{green!60!black},
  morecomment=[l][\color{magenta}]{\#},
  numbers=left,
  numberstyle=\tiny\color{gray},
  stepnumber=1,
  numbersep=10pt,
  backgroundcolor=\color{gray!10},
  frame=single,
  breaklines=true,
  tabsize=2,
  captionpos=b,
    literate=
  {ç}{{\c c}}1
  {Ç}{{\c C}}1
  {ã}{{\~a}}1
  {õ}{{\~o}}1
  {á}{{\'a}}1
  {à}{{\`a}}1
  {é}{{\'e}}1
  {í}{{\'i}}1
  {ó}{{\'o}}1
  {ú}{{\'u}}1
}
% ---
		
% ---
% Pacotes adicionais, usados apenas no âmbito do Modelo Canônico do abnteX2
% ---
\usepackage{lipsum}				% para geração de dummy text
% ---
		
% ---
% Pacotes de citações
% ---
\usepackage[brazilian,hyperpageref]{backref}	 % Paginas com as citações na bibl
\usepackage[alf]{abntex2cite}	% Citações padrão ABNT
% ---

% ---
% Configurações do pacote backref
% Usado sem a opção hyperpageref de backref
\renewcommand{\backrefpagesname}{Citado na(s) página(s):~}
% Texto padrão antes do número das páginas
\renewcommand{\backref}{}
% Define os textos da citação
\renewcommand*{\backrefalt}[4]{
	\ifcase #1 %
		Nenhuma citação no texto.%
	\or
		Citado na página #2.%
	\else
		Citado #1 vezes nas páginas #2.%
	\fi}%
% ---

% ---
% Informações de dados para CAPA e FOLHA DE ROSTO
% ---
\titulo{Comparação de Eficiência entre\\ Algoritmos de Ordenação}
\autor{ Ana Beatriz Vasconcelos \and João Monteiro}
\local{Brasil}
\data{2026}
% ---

% ---
% Configurações de aparência do PDF final

% alterando o aspecto da cor azul
\definecolor{blue}{RGB}{41,5,195}

% informações do PDF
\makeatletter
\hypersetup{
     	%pagebackref=true,
		pdftitle={\@title}, 
		pdfauthor={\@author},
    	pdfsubject={Modelo de artigo científico com abnTeX2},
	    pdfcreator={LaTeX with abnTeX2},
		pdfkeywords={abnt}{latex}{abntex}{abntex2}{atigo científico}, 
		colorlinks=true,       		% false: boxed links; true: colored links
    	linkcolor=blue,          	% color of internal links
    	citecolor=blue,        		% color of links to bibliography
    	filecolor=magenta,      		% color of file links
		urlcolor=blue,
		bookmarksdepth=4
}
\makeatother
% --- 

% ---
% compila o indice
% ---
\makeindex
% ---

% ---
% Altera as margens padrões
% ---
\setlrmarginsandblock{3cm}{3cm}{*}
\setulmarginsandblock{3cm}{3cm}{*}
\checkandfixthelayout
% ---

% --- 
% Espaçamentos entre linhas e parágrafos 
% --- 

% O tamanho do parágrafo é dado por:
\setlength{\parindent}{1.3cm}

% Controle do espaçamento entre um parágrafo e outro:
\setlength{\parskip}{0.2cm}  % tente também \onelineskip

% Espaçamento simples
\SingleSpacing

% ----
% Início do documento
% ----
\begin{document}

% Retira espaço extra obsoleto entre as frases.
\frenchspacing 

% ----------------------------------------------------------
% ELEMENTOS PRÉ-TEXTUAIS
% ----------------------------------------------------------

%---

%
%---
% página de titulo
\maketitle

% resumo em português
\begin{resumoumacoluna}
inserir resumo
 \vspace{\onelineskip}
 
 \noindent
 \textbf{Palavras-chaves}: Métodos de Ordenação. Estrutura de Dados e Algoritmos. Análise de Eficiência.
\end{resumoumacoluna}



% ----------------------------------------------------------
% ELEMENTOS TEXTUAIS
% ----------------------------------------------------------
\textual

% ----------------------------------------------------------
% Introdução
% ----------------------------------------------------------
\section*{Introdução}
\addcontentsline{toc}{section}{Introdução}
Algoritmos podem ser entendido como uma série de passos para transformar um input em um output desejado. Segundo \cite{cormen2022introduction} um algoritmo pode ser entendido informalmente como um procedimento computacional bem definido que recebe um valor ou um conjunto de valores, como inputs e produz um valor ou conjutno de valores, como outputs.

Um tópico frequente e bastante útil na programação é o ordenamento de itens e de acordo com \cite{donald1999art} existem diversas aplicações de algoritmos de ordenação e busca dentre elas: busca por informação por chave e valor, matching items em um ou mais arquivos e resolver o problema de"\textit{togetherness}". Para resolver problemas de ordenação existem diferentes algoritmos que resolvem o problema de forma diversa, sendo a compreensão e o estudo do suas particularidades fundamental para selecionar o algoritmo que melhor se aplica a problemas praticos de ordenação.

A fim de avaliar o desempenho em termos de tempo de execução dentre os diferentes algoritmos de ordenção, neste trabalho realizaremos uma análise empirica a partir da comparação do tempo de execução de dos principais algoritmos de ordenação sendo eles: \textit{buble sort, insertion sort, selection sort, merge sort, quick sort}.

% ----------------------------------------------------------
% Fundamentação teórica
% ----------------------------------------------------------
\section{Fundamentação teórica}


\subsection{Bubble sort}
O método bubble sort ou bolha consiste em comparar os elementos dois a dois. Em cada iteração desse método a lista é percorrida a partir do seu início comparando cada elemento com seu sucessor, trocando-os de posição caso se aplique. No método bolha os elementos mais "densos" são passados para o "fundo" da lista, ou seja, para uma lista de tamanho n, um elemento que está inicialmente numa posição i e, para obter uma lista ordenada, ele deve ocupar a posição j, então ele terá que passar por todas as posições entre i e j.\cite{viana2011pesquisa}.

A Figura \ref{fig:bubblesort} apresenta a implementação do método bolha que consiste em dois laços de repetição sendo o mais externo o que determinará o total de iterações a ser realizada, enquanto o laco interno percorre a lista comparando os elementos e realizando as trocas iniciando em 0 até n - 1 - i, ou seja, a cada iteração o método estará ordenando uma lista menor uma vez que os items nas ultimas posições i do vetor já estão ordenados. Adicionalmente, é utilizada uma variável booleana \textit{swapped}, que permite identificar se houve alguma troca durante uma iteração. Caso nenhuma troca ocorra, conclui-se que a lista já está ordenada, permitindo a interrupção antecipada do algoritmo e evitando execuções desnecessárias.
\begin{figure}[H]
\centering
\begin{lstlisting}[]
class BubbleSort {
public:
  BubbleSort() {};
  ~BubbleSort() {};

  void sort(std::vector<int> &arr) {
    std::cout << "Iniciando Bubble Sort..." << std::endl;
    int n = arr.size();
    bool swapped;

    for (int i = 0; i < n - 1; i++) {
      swapped = false;

      for (int j = 0; j < n - 1 - i; j++) {
        if (arr[j] > arr[j + 1]) {
          std::cout << "Trocando " << arr[j] << " e " << arr[j + 1] << endl;
          std::swap(arr[j], arr[j + 1]);
          swapped = true;
        }
      }

      if (!swapped) {
        break;
      }
    }
  }
};
\end{lstlisting}
\caption{Implementação do algoritmo Bubble Sort em C++}
\label{fig:bubblesort}
\end{figure}

No método bolha a operação de troca e comparação são as operações relevantes e analisando a complexidade no pior caso desse algoritmo temos:

\begin{equation}
\begin{gathered}
t(1) = n - 1 \\
t(2) = n - 2 \\
t(3) = n - 3 \\
\vdots \\
t(n-1) = 1 \\
\\
T(n) = (n-1) + (n-2) + \cdots + 1 \\
\\
T(n) = \sum_{i=1}^{n-1} i \\
\end{gathered}
\end{equation}
Aplicando a fórmula da soma de uma progressão aritmética (PA), temos:
\[
\frac{(n-1 + 1)\,n}{2}.
\]

Simplificando, obtemos:
\[
\frac{n(n-1)}{2}.
\]

Baseada em técnicas apresentadas em material de apoio da disciplina de estrutura de dados \cite{bigO2026} podemos aplicar a notação Big O:
\[
O(n^2 - n).
\]

Priorizando o termo de maior grau, é possivel concluir que a complexidade do método bolha no pior caso é:
\[
O(n^2).
\]


\subsection{Insertion sort}
O método de insertion sort desloca vários elementos e insere um elemento na ordem correta. Nesse metodo cada interação consiste em inserir o elemento na sua ordem correta, percorrendo a lista da esquerda para direita. A Figura \ref{fig:insertionsort} apresenta a implementação do algoritmo de inserção. Primeiro obtem-se o número de elementos do vetor (n) então o laço de repetição percorre o vetor de 1 até n definindo a variavel numérica \textit{key} que é utilizada para realizar as comparações e j é a posição anterior a i. No laço de repetição interno verifica-se se j é maior ou igual a zero, ou seja, ainda não se chegou ao fim do vetor e se j é maior que a \textit{key} utilizada. Enquanto essa condição é valida, o elemento na posição j será movido para j + 1 e j é reduzido em uma unidade. Por fim, o elemento de interesse, \textit{key}, é introduzido na posição j + 1. 


\begin{figure}[H]
\centering
\begin{lstlisting}[language=C++]
class InsertionSort {
public:
    InsertionSort() {};
    ~InsertionSort() {};

    void sort(std::vector<int> &arr){
        std::cout << "Iniciando Insertion Sort..." << std::endl;
        int n = arr.size();

        for (int i = 1; i < n; i++){
            int key = arr[i];
            int j = i - 1;

            // Enquanto o elemento à esquerda for maior que a chave
            // empurramos ele uma posicão para a direita
            while(j >= 0 && arr[j] > key){
                std::cout << "Movendo " << arr[j] << " para a direita" << std::endl;
                arr[j + 1] = arr[j];
                j = j - 1;
            }

            arr[j + 1] = key;
        }
    }
};
\end{lstlisting}
\caption{Implementação do algoritmo Insertion Sort em C++}
\label{fig:insertionsort}
\end{figure}


Ao analisar o pior caso, uma lista inversamente ordenada temos:
Para o laço externo:
\[
T(n) = (n - 1) \\
\]
Para o laço de interno:
\begin{equation}
\begin{gathered}
t(1) = 1 \\
t(2) = 2 \\
t(3) = 3 \\
\vdots \\
t(n-1) = n -1 \\
\\
T(n) = (1) + (2) + \cdots + (n - 1) \\
\\
T(n) = \sum_{i=1}^{n-1} i \\
\end{gathered}
\end{equation}
Aplicando a fórmula da soma de uma progressão aritmética (PA), temos:
\[
\frac{(1 + (n - 1))\,(n -1)}{2}.
\]

Simplificando, obtemos:
\[
\frac{n(n-1)}{2}.
\]

Em notação Big O:
\[
O(n^2 - n).
\]

Priorizando o termo de maior grau, é possivel concluir que a complexidade do método de inserção no pior caso é:
\[
O(n^2).
\]
\subsection{Selection Sort}
 O metodo de seleção ou \textit{selection sort} assim como o de inserção insere o elemento na posição correta a partir da seleção do elemento de maior prioridade. a figura \ref{fig:selectionsort} apresenta um exemplo de implementação desse metodo. 

\begin{figure}[H]
\centering
\begin{lstlisting}[]
class SelectionSort {
public:
    SelectionSort() {}
    ~SelectionSort() {}

    void sort(std::vector<int> &arr) {
        std::cout << "Iniciando Selection Sort..." << std::endl;
        int n = arr.size();

        for (int i = 0; i < n - 1; i++) {
            int min_idx = i;
            for (int j = i + 1; j < n; j++) {
                if (arr[j] < arr[min_idx]) {
                    min_idx = j;
                }
            }

            if (min_idx != i) {
                std::cout << "Trocando " << arr[i] 
                          << " pelo menor encontrado " 
                          << arr[min_idx] << std::endl;
                std::swap(arr[i], arr[min_idx]);
            }
        }
    }
};
\end{lstlisting}
\caption{Implementação do Selection Sort em C++}
\label{fig:selectionsort}
\end{figure}

Ao analisar o método é possivel destacar como operações relevantes as operações de comparação e troca o que permite a análise complexidade do algoritmo sendo:
para cada execução 
\begin{equation}
\begin{gathered}
t(2) =  1 \\
t(3) = 3 \\
t(4) = 6 \\
t(5) = 10 \\
\vdots \\
t(n-1) = n -1  \\
\\
T(n) = 1 + 2 + \cdots + (n - 1) \\
\\
T(n) = \sum_{i=1}^{n-1} i \\
\end{gathered}
\end{equation}
Aplicando a fórmula da soma de uma progressão aritmética (PA), temos:
\[
\frac{(n-1 + 1)\,n}{2}.
\]

Simplificando, obtemos:
\[
\frac{n(n-1)}{2}.
\]

Ao aplicar a notação Big O:
\[
O(n^2 - n).
\]

Priorizando o termo de maior grau, é possivel concluir que a complexidade do método de seleção no pior caso é:
\[
O(n^2).
\]

\subsection{Merge Sort}

O \textit{merge sort} soluciona o problema de ordenação a partir da divisão do vetor em duas partes de igual tamanho e ordena os subconjuntos para depois realizar a mesclagem ou \textit{merge} dos subconjuntos. A figura \ref{fig:mergesort} apresenta um exemplo de implementação do método de mesclagem.

\begin{figure}[H]
\centering
\begin{lstlisting}[]
class MergeSort {
public:
    MergeSort() {}
    ~MergeSort() {}

    void merge(std::vector<int> &arr, int left, int middle, int right) {
        int n1 = middle - left + 1;
        int n2 = right - middle;

        std::vector<int> L(n1), R(n2);

        for (int i = 0; i < n1; i++) {
            L[i] = arr[left + i];
        }
        for (int j = 0; j < n2; j++) {
            R[j] = arr[middle + 1 + j];
        }

        int i = 0, j = 0, k = left;

        while (i < n1 && j < n2) {
            if (L[i] <= R[j]) {
                arr[k] = L[i];
                i++;
            } else {
                arr[k] = R[j];
                j++;
            }
            k++;
        }

        while (i < n1) arr[k++] = L[i++];
        while (j < n2) arr[k++] = R[j++];

        std::cout << "Mesclando sub-arrays de "
                  << left << " a " << right << std::endl;
    }

    void sortRecursive(std::vector<int> &arr, int left, int right) {
        if (left < right) {
            int middle = left + (right - left) / 2;

            sortRecursive(arr, left, middle);
            sortRecursive(arr, middle + 1, right);
            merge(arr, left, middle, right);
        }
    }

    void sort(std::vector<int> &arr) {
        std::cout << "Iniciando Merge Sort..." << std::endl;
        if (arr.size() > 1) {
            sortRecursive(arr, 0, arr.size() - 1);
        }
    }
};
\end{lstlisting}
\caption{Implementação do Merge Sort em C++}
\label{fig:mergesort}
\end{figure}
% ----------------------------------------------------------
% Metodologia
% ----------------------------------------------------------
% ---

\section{Metodologia}
Foram elaborados testes para cada um dos algoritmos de ordenação estudados a partir de construção de um programa e de vetores de diferentes tamanhos iniciando de 5000 ate 20000, aumentando de 1000 em 1000 até obter o vetor de maior tamanho. A fim de analisar o tempo médio de execução de cada algoritmo estudado para cada tamanho de vetor, foram realizadas multiplas execuções com vetores preenchidos a partir de números aleatórios em um intervalo de 0 até 100000.

\subsection{Projeto/Programa}

Descrever organização do projeto /libs etc 

% Finaliza a parte no bookmark do PDF, para que se inicie o bookmark na raiz
% ---
\bookmarksetup{startatroot}% 
% ---

% ---
% Conclusão
% ---
\section*{Considerações finais}
\addcontentsline{toc}{section}{Considerações finais}

\lipsum[1]

\begin{citacao}
\lipsum[2]
\end{citacao}

\lipsum[3]

% ----------------------------------------------------------
% ELEMENTOS PÓS-TEXTUAIS
% ----------------------------------------------------------
\postextual


% ----------------------------------------------------------
% Referências bibliográficas
% ----------------------------------------------------------
\bibliography{abntex2-modelo-references}

\end{document}
